. Key properties that make LLMs useful in the RL setting
include: (a) \emph{emergent instruction-following}, (b) \emph{in-context
learning} from demonstrations and feedback provided in the prompt,
(c) \emph{code generation}, enabling translation from natural language to
executable programs, and (d) \emph{chain-of-thought reasoning}, enabling
multi-step deliberation.


 Three trends have driven the convergence. First, \emph{alignment}: as LLMs
are deployed in consequential settings, the need to shape their long-horizon
behavior resembles the RL objective---optimize for some notion of ``good''
behavior over a sequence of decisions. Second, \emph{agency}: LLMs are
increasingly deployed as agents that act in environments (code executors, web
browsers, robotic controllers), making the RL framework natural. Third,

=======≠≠≠≠≠≠=========≠≠≠=
1) structure verbal feedback into actionable indights. 
2) design agents which takes verbal feedback
3) investment in tools; redesign tools to
4) communication protocol where tool should be able to communicate with LLM in a way where they understand each other. 
5) Latency for verbal feedback

==============
